\section{最小混合試験族の三次中心元の閉公式}
\label{mt:sec:family}

一楕円ブロックの分離型から一般混合型へ進むため、
五変数のcubic
\begin{equation}\label{mt:eq:classical}
 F_\lambda=\frac{x^3+y^3+z^3}{3}+u^3+v^3+\lambda uxy
\end{equation}
を固定する。$u,v$を中心、$h,\lambda$を次数零のパラメータとし、
\begin{equation}\label{mt:eq:relations}
 [y,z]=h(x^2+\lambda uy),\qquad
 [z,x]=h(y^2+\lambda ux),\qquad
 [x,y]=hz^2
\end{equation}
で周囲の代数$A_{h,\lambda}$を定める。
本節ではこの族の三次中心元を多項式の閉公式で構成し、
自由代数内の三つの恒等式によって中心性を証明する。
これは可換復元先の二次係数$G_2$の計算ではなく、
それを計算するための非可換代数側の入力である。
中心式の特殊形は既存の混合中心式の有限体検算に現れるが、
ここでは特殊化を用いず、全パラメータに対する証明書を与える。
閉公式の先行研究上の初出は未確認である。

\begin{lemma}[古典特殊ファイバーの滑らかさ]
\label{mt:lem:classical-smooth}
$V(F_\lambda)\subset\PP^4$は$\lambda^3\ne-3$のとき、かつそのときに限り滑らかである。
\end{lemma}
\begin{proof}
全偏微分が零なら$z=v=0$であり、残る方程式は
\[
 x^2+\lambda uy=y^2+\lambda ux=3u^2+\lambda xy=0.
\]
$u=0$なら$x=y=0$となるので射影点は存在しない。
$u=1$とすると、$\lambda=0$では解がなく、$\lambda\ne0$なら$x,y$は非零である。
最初の二式を掛けると$xy=\lambda^2$を得て、第三式は$3+\lambda^3=0$となる。
逆にこの等式が成立すれば、$x^3=3$を満たす$x$と
$y=-x^2/\lambda$を選ぶことで解を得る。
三次のEuler等式により、これらの解は$F_\lambda=0$も満たす。
\end{proof}

\subsection{中心三次式と自由代数の証明書}

$a=\lambda u$と略記し、$T=xyz+xzy+yxz+yzx+zxy+zyx$と置く。

\begin{proposition}[混合試験族の中心元]
\label{mt:prop:center}
\begin{equation}\label{mt:eq:center}
 \boxed{\begin{aligned}
 \Omega_{h,\lambda}={}&
 \frac{x^3+y^3+z^3}{3}-\frac{h^2}{9}T
 +\frac a2(xy+yx)\\
 &+\frac{h^2a^2}{6}z+\frac{h^4a}{18}z^2+u^3+v^3.
 \end{aligned}}
\end{equation}
は$A_{h,\lambda}$の三次中心元である。
これは$\CC[h,\lambda]$上の恒等的な中心性であり、
$h\ne0$への局所化も、有限$h$-jetへの制限も必要としない。
$h=0$で$\Omega_{0,\lambda}=F_\lambda$である。
\end{proposition}
\begin{proof}
まず$a$を独立な中心係数とし、自由代数
$\mathbb Q[h,a]\langle x,y,z\rangle$で計算する。
関係式を
\[
 r_1=xy-yx-hz^2,\quad
 r_2=xz-zx+hy^2+hax,\quad
 r_3=yz-zy-hx^2-hay
\]
と書く。二つの演算を
\[
 \mathcal S_p(r)=rp^2+prp+p^2r,\qquad
 \mathcal D_{p,q}(r)=rpq-prq-qrp+qpr
\]
とし、$\{p,q\}_+=pq+qp$とする。
$\omega=\Omega_{h,\lambda}-u^3-v^3$について、以下の三式が
\emph{関係式で割る前の自由代数}において成立する：
\begin{align}
 [\omega,x]={}&
 -\frac13\mathcal S_y(r_1)-\frac13\mathcal S_z(r_2)
 +\frac{h^2}{18}
   \bigl(\mathcal D_{z,x}(r_1)+\mathcal D_{y,x}(r_2)\bigr)\notag\\
 &+\frac{ah^3}{18}[x,r_1]-\frac a2\{x,r_1\}_+
  +\frac{ah}{6}[z,r_2]+
  \frac h6[yz+zy,r_3],\label{mt:eq:certificate-x}\\
 [\omega,y]={}&
 \frac13\mathcal S_x(r_1)-\frac13\mathcal S_z(r_3)
 -\frac{h^2}{9}\mathcal D_{z,y}(r_1)
 +\frac{ah^3}{18}[y,r_1]+\frac a2\{y,r_1\}_+\notag\\
 &+\frac{ah^2}{18}[r_2,x]+\frac h9[r_2,2xz+zx]
 +\frac{h^2}{9}[r_2,y]y+\frac{ah}{6}[r_3,z],
 \label{mt:eq:certificate-y}\\
 [\omega,z]={}&
 \frac{a(2h^3+9)}{18}[r_1,z]
 +\frac h9[2xy+yx,r_1]+\frac{h^2}{9}[r_1,z]z\notag\\
 &+ayr_2+\frac13\mathcal S_x(r_2)
 -\frac{h^2}{9}\mathcal D_{y,z}(r_2)
 +ar_3x+\frac13\mathcal S_y(r_3).
 \label{mt:eq:certificate-z}
\end{align}
各式は語と$h,a$の冪ごとに展開して確認でき、全係数は有理数である。
右辺の各項は$r_i$を含むので、商代数で全交換子が零となる。
$a=\lambda u$を代入し、中心スカラー$u^3+v^3$を戻せば中心性を得る。
零パラメータでの表示は直接代入による。
\end{proof}

\subsection{Artin切断でのsorted PBW表示}

閉公式の中心性と、周囲の代数のflatnessを別々に確認する。
$u,v$の中心性を先に使い、active語を$z<y<x$の順に並べる。
還元規則は
\begin{equation}\label{mt:eq:rewrite}
 xy\longmapsto yx+hz^2,\qquad
 xz\longmapsto zx-hy^2-hax,\qquad
 yz\longmapsto zy+hx^2+hay
\end{equation}
である。これは一般の非零$h$で通常の語順序に関して停止する規則とは主張しない。

\begin{lemma}[全Artin切断でのPBW性]
\label{mt:lem:artin-pbw}
任意の$N\ge0$について、$h^{N+1}=0$を課した代数は
$\CC[\lambda,u,v,h]/(h^{N+1})$上で
$\{z^iy^jx^k:i,j,k\ge0\}$を基底に持つ。
従って各次数成分は$\CC[\lambda,h]/(h^{N+1})$上自由であり、
周囲のHilbert級数は$(1-t)^{-5}$である。
\end{lemma}
\begin{proof}
$h^r w$の還元に、辞書式に並べた二つの非負整数
$(N-r,\operatorname{inv}(w))$を用いる。
$h$を含まない交換項では反転数が減り、補正項では$h$の冪が増える。
従って$h^{N+1}=0$のもとで全還元が停止する。
active語の重なりは$xyz$の一つだけである。
最初に$xy$を還元した場合と、最初に$yz$を還元した場合は、いずれも有限回の置換で
\begin{equation}\label{mt:eq:overlap}
 zyx+h(x^3-y^3+z^3)
\end{equation}
になる。$a$を含む項は両経路でそれぞれ相殺する。
包含の曖昧さはない。係数に逆元を導入しないmonicな合流的還元なので、
既約語$z^iy^jx^k$が基底となる。
中心係数$a$を$\lambda u$へ置き換え、独立な中心変数$v$を加えても同じ証明が成り立つ。
\end{proof}

この補題は全ての$h$-Artin切断およびその逆極限についてのformal PBW性を示す。
非零$h$での全ファイバーのAS-regular性や代数的族全体のflatnessを、
Artin切断の計算だけから結論していない。
また、中心商は$v$に関してmonicな三次式であるため、
各Artin切断で$1,v,v^2$を、$v$を含まない周囲の代数上の基底に持つ。
従って中心三次商も次数別に平坦で、Hilbert級数は
$(1-t^3)/(1-t)^5$である。

